% AUTO-GENERATED by scripts/generate_paper_results.py from
% evaluation/controlled/results.json. Do not hand-edit.
\begin{table}[t]
\centering
\small
\begin{tabular}{lrrr}
\toprule
\textbf{Category} & \textbf{Recall} & \textbf{Acc.\ given retrieved} & \textbf{N} \\
\midrule
direct       & 50.0\% & 100.0\% & 2 \\
negated      & 100.0\% & 100.0\% & 2 \\
numeric      & 100.0\% & 100.0\% & 2 \\
temporal     & 50.0\% & 100.0\% & 2 \\
entity       & 100.0\% & 100.0\% & 2 \\
attribute    & 100.0\% & 100.0\% & 2 \\
\midrule
\emph{all CONTRADICTS} & 83.3\% & 100.0\% & 12 \\
SUPPORTS & 66.7\% & 25.0\% & 6 \\
NEUTRAL & 0.0\% & -- & 6 \\
\midrule
\textbf{Overall} & \textbf{58.3\%} & \textbf{78.6\%} & \textbf{24} \\
\bottomrule
\end{tabular}
\caption{SMRITI-Controlled-v1 results, post-fix. Retrieval recall 95\% CI
(bootstrap, $n = 24$): [37.5\%, 79.2\%].
Classification accuracy given retrieval, 95\% CI ($n = 14$): [57.1\%, 100.0\%].}
\label{tab:controlled}
\end{table}


Table~\ref{tab:controlled} is, we think, the cleanest evidence in this
paper. Three findings, each requiring no interpretation at all:

\textbf{(1) Across all six contradiction categories, every retrieved pair
was classified correctly}: 10 of 12 designed contradiction pairs were
retrieved as NLI candidates, and all 10 were correctly labeled
CONTRADICTS with contradiction scores of 0.996--1.000 --- 100\%
classification precision on a fully objective test set, spanning direct
reversal, explicit negation, numeric facts, temporal state, named-entity
attribution, and property negation alike. This directly answers RQ2/RQ3
for \emph{propositional} contradiction: the arg-max and relatedness-gate
fix is not an artifact of one corpus's idiosyncrasies. It is not the
final word on RQ3, however --- \S\ref{sec:results-annotation} shows the
same fix does not transfer to SMRITI-Reference's \emph{rhetorical}
disagreement, which this controlled corpus was not designed to test.

\textbf{(2) All six ``easy'' NEUTRAL control pairs were correctly
excluded before NLI ever ran} (0/6 retrieved as candidates at all) ---
i.e.\ zero false-positive contradictions on pairs sharing no vocabulary
with their partner, confirming that embedding-based candidate retrieval
correctly filters clearly-unrelated pairs. This is a weaker claim than
it may look: it shows retrieval works, not that the relatedness gate
would have rejected a false CONTRADICTS had one of these pairs actually
been retrieved (an external review point) --- \S\ref{sec:hard-neutral}
tests that directly.

\textbf{(3) SUPPORTS is the weaker link}: only 4 of 6 designed SUPPORTS
pairs were retrieved, and of those, only 1 was correctly labeled SUPPORTS
(the other 3 were split across NEUTRAL and REFINES). REFINES is, in one
case, a defensible alternative reading ("X improves Y" followed by "an
independent benchmark confirms X improves Y" is arguably an elaboration as
much as a confirmation), but the low classification accuracy here --- in
sharp contrast to CONTRADICTS's 100\% --- indicates the entailment side of
the resolver (which received the identical arg-max fix as the
contradiction side) is not yet as well-calibrated, and is a concrete
target for future work distinct from the contradiction fix this paper
centers on.

\subsection{Hard-Neutral Extension: Does the Relatedness Gate Actually
Work?}
\label{sec:hard-neutral}

% AUTO-GENERATED by scripts/generate_paper_results.py from
% evaluation/controlled/results.json. Do not hand-edit.
\begin{table}[t]
\centering
\small
\begin{tabular}{lrrrr}
\toprule
\textbf{Hard-neutral category} & \textbf{N} & \textbf{Retrieved} & \textbf{Recall} & \textbf{Misclassified} \\
\midrule
Same topic, diff.\ proposition & 2 & 1 & 50.0\% & 1 \\
Lexical collision & 2 & 0 & 0.0\% & 0 \\
Adjacent, non-conflicting & 2 & 1 & 50.0\% & 0 \\
\bottomrule
\end{tabular}
\caption{SMRITI-Controlled's hard-neutral extension: NEUTRAL pairs
deliberately constructed to be topically related (unlike the original
six ``easy'' NEUTRAL pairs, which share no vocabulary at all and are
never retrieved). ``Misclassified'' counts retrieved pairs the resolver
labeled anything other than NEUTRAL --- the false-CONTRADICTS/SUPPORTS/
REFINES rate on pairs the relatedness gate cannot reject by construction,
since they genuinely share topic or vocabulary.}
\label{tab:hard-neutral}
\end{table}


\textbf{(4) The relatedness gate has a real, demonstrated gap.} Of the
six hard-neutral pairs, two were retrieved as NLI candidates at all
(the other four --- both lexical-collision pairs and one each of the
other two categories --- were not retrieved, so the relatedness gate was
never exercised on them). Of those two, \textbf{one was misclassified
CONTRADICTS}: ``The Aerotrix rover uses LiDAR for obstacle detection''
paired with ``The Aerotrix rover uses a stereo camera for obstacle
detection'' --- two claims about the same named entity, sharing the
entity token ``Aerotrix'' (exactly the token overlap the relatedness
gate's entity-token bonus is designed to detect as \emph{related}), but
describing two compatible sensors, not a genuine conflict. The
relatedness gate correctly identified the pair as related --- that is
its entire job --- and passed it through to the resolver, which then
read ``uses LiDAR'' and ``uses a stereo camera'' as mutually exclusive
rather than as two attributes a single rover could have. This is
precisely the gap an external review predicted: the gate only asks
\emph{is this pair topically related}, never \emph{are these two
specific attributes of this entity actually incompatible}, and the NLI
model has no special handling for multi-attribute entities either. The
other retrieved hard-neutral pair (the two slow-braising statements,
adjacent-non-conflicting category) was correctly labeled NEUTRAL,
so this is not a uniform failure of the gate; with $n=1$ misclassified
pair, we report this as a demonstrated failure mode, not a rate, and
record it as E7 in \S\ref{sec:error-taxonomy}.
